\section{All circular convolution powers of log Gamma}
\label{stconv:sec:gamma}
At the integrated end no finite part is needed. Put
\begin{equation}\label{stconv:eq:gr}
 L=\log(2\pi),\qquad Z_j(x)=\zeta^{[j]}(0,x),\qquad
 g(x)=Z_1(x)=\log\Gamma(x)-\frac L2.
\end{equation}
Each $Z_j$ has mean zero. It is locally a polynomial in $\log x$ plus
a smooth function near zero, and hence belongs to every finite $L^p(0,1)$.

\begin{theorem}[All mixed integrated jets]\label{stconv:thm:integrated}
Let $d\ge2$ and $r_1,\ldots,r_d\ge0$ be integers. Write
$S=s_1+\cdots+s_d$. Then, as an identity of continuous functions after
circular convolution,
\begin{equation}\label{stconv:eq:mixedintegrated}
 Z_{r_1}*\cdots*Z_{r_d}(x)
 =\left.
 \partial_{s_1}^{r_1}\cdots\partial_{s_d}^{r_d}
 \left[\frac{\prod_{i=1}^d\Gamma(1-s_i)}{\Gamma(d-S)}
              \zeta(1-d+S,x)\right]\right|_{s_1=\cdots=s_d=0}.
\end{equation}
It holds for $0<x<1$ and at the integer point by continuous extension.
Thus any mixed convolution of the centered primitives $A_n$ reduces to
finitely many jets at the single integer $1-d$.
\end{theorem}
\begin{proof}
Before differentiation, write each Hurwitz factor as
$\zeta(s_i,x)=\Gamma(1-s_i)R_{1-s_i}(x)$ and apply
Theorem~\ref{stconv:thm:semigroup} repeatedly. For parameters sufficiently close
to zero, each factor is holomorphic as an $L^2$-valued family: the
endpoint singularities are bounded by
$x^{-\eta}(1+|\log x|^M)$ with $2\eta<1$ on each fixed derivative
order and compact parameter set. Convolution of two $L^2$ functions is
continuous, by Cauchy--Schwarz and translation continuity; further
convolution with $L^1$ factors preserves continuity. These same estimates
justify differentiating all the convolutions. The right side is
continuous at the integer point since $\Rea(1-d+S)<0$ in a small
parameter neighborhood. This proves the formula and its endpoint
extension.
\end{proof}

The case with every $r_i=1$ admits a compact all-$d$ formula. Define
$H_m^{(j)}=\sum_{a=1}^ma^{-j}$, and, for fixed $d\ge2$, set
\begin{equation}\label{stconv:eq:bd}
 b_1=H_{d-1},\qquad
 b_j=(j-1)!\big[H_{d-1}^{(j)}-\zeta(j)\big]\quad(j\ge2).
\end{equation}

\begin{theorem}[All convolution powers of centered log-gamma]\label{stconv:thm:gammapowers}
For every integer $d\ge2$,
\begin{equation}\label{stconv:eq:gammapowers}
 g^{*d}(x)=\frac1{(d-1)!}\sum_{j=0}^d\binom dj
       \stconvBell_{d-j}(b_1,\ldots,b_{d-j})\zeta^{[j]}(1-d,x),
\end{equation}
where $\stconvBell_0=1$ and the $b_j$ are scalar quantities in~\eqref{stconv:eq:bd}.
\end{theorem}
\begin{proof}
In~\eqref{stconv:eq:mixedintegrated}, differentiate once in each $s_i$.
Only the squarefree coefficient $s_1\cdots s_d$ matters. Modulo
$s_i^2=0$, one has
\[
 \prod_{i=1}^d\Gamma(1-s_i)
 =\prod_{i=1}^d(1+\gamma s_i)
 =\exp(\gamma S).
\]
Consequently the required mixed derivative equals the ordinary $d$th
derivative at $S=0$ of
$e^{\gamma S}\zeta(1-d+S,x)/\Gamma(d-S)$. Now
\begin{align*}
 \log\frac{e^{\gamma S}\Gamma(d)}{\Gamma(d-S)}
 &=H_{d-1}S+
   \sum_{j\ge2}\frac{H_{d-1}^{(j)}-\zeta(j)}jS^j.
\end{align*}
This follows directly from
$\Gamma(d-S)=\Gamma(1-S)\prod_{a=1}^{d-1}(a-S)$ and the Taylor series
of $\log\Gamma(1-S)$. Apply~\eqref{stconv:eq:Belldef} and the Leibniz rule;
$\Gamma(d)=(d-1)!$ supplies the prefactor.
\end{proof}

The first two nontrivial cases are
\begin{align}
 g*g={}&\zeta''(-1,x)+2\zeta'(-1,x)
                   +[2-\zeta(2)]\zeta(-1,x),\label{stconv:eq:g2}\\
 g^{*3}={}&\frac12\zeta'''(-2,x)+\frac94\zeta''(-2,x)
       +\left[\frac{21}4-\frac32\zeta(2)\right]\zeta'(-2,x)
       \label{stconv:eq:g3}\\
 &+\left[\frac{45}8-\frac94\zeta(2)-\zeta(3)\right]\zeta(-2,x).
 \notag
\end{align}
The first left side is explicitly
\[
 \int_0^xg(t)g(x-t)\dd t+\int_x^1g(t)g(1+x-t)\dd t.
\]
For uncentered $f=\log\Gamma$, mean-zero convolution annihilates all
mixed constant/nonconstant terms, so
\begin{equation}\label{stconv:eq:uncentered}
 f^{*d}(x)=\left(\frac L2\right)^d+g^{*d}(x).
\end{equation}

\begin{remark}[Convolution powers are not pointwise moments]
The value $g^{*3}(0)$ is a two-dimensional circular-convolution integral.
It is not $\int_0^1g(t)^3\dd t$. Thus~\eqref{stconv:eq:g3} does not eliminate
the Tornheim derivative coordinates of the ordinary cubic log-gamma
moment discussed in the repository. Likewise, the mixed reflected
pointwise moments in \texttt{07-reflected-moments.tex} are different
objects. This distinction prevents a spurious claim of resolving that
separate reduction problem.
\end{remark}
